\scanpage{101}
% archive foliation: 100
On s'intéresse surtout au schéma $(M_{\mathcal C_0})_{\mathbb Q}$ --- qui est donc fini étale sur $\mathbb Q$, non ramifié en les $p>\delta$.

\deleted{\unsure{Sur ce schéma}}
\deleted{Je \unsure{le désignerai}}

Sur le schéma $M_{\mathcal C_0}$, le gr. $\gamma_{\mathcal C_0}=\operatorname{Aut}(\mathcal C_0)$ opère de façon évidente car il opère sur $F_{\mathcal C_0}$. Le stabilisateur d'une section de $M_{\mathcal C_0}$ sur un $T$, i.e. d'une classe d'iso de ) $X\xrightarrow{f}\mathbb P^1_S$ \unsure{\ldots} avec \deleted{$\unsure{\gamma}$} $\mathcal C_0$-\deleted{\unsure{ori}}épinglage, correspond \inserted{(si $S$ est connexe)} aux autom du \deleted{$X\unsure{-f}$} rev. $X\xrightarrow{f}\mathbb P^1_S$ indépendamment de l'épinglage. P. ex. si $f$ est galoisien, il est can. iso. au groupe de Galois de $f$.

Pour.\ $\unsure{\forall}$ classe d'iso. $\alpha$ de \deleted{cartes}\inserted{groupes} \deleted{(non épinglées)} \inserted{finis}, \deleted{orientées (de type $\mathcal C_0$)} \inserted{\unsure{(constant fini)}} soit $F_{\mathcal C_0}^{\alpha}$ le sous-foncteur \deleted{que} $F_{\mathcal C_0}$ qui correspond aux $X\xrightarrow{f}\mathbb P^1_T$ tels que sur \inserted{chaque} fibre géom., le groupe de Galois du
